(1) 以第 \(n\) 个重原子 \(u_n\) 与轻原子 \(v_n\) 为一组，连接方式取重原子同时接 \(v_n\)、\(v_{n-1}\)，轻原子接 \(u_n,u_{n+1}\)。牛顿方程 \(M\ddot u_n=K(v_n+v_{n-1}-2u_n)\)，\(m\ddot v_n=K(u_n+u_{n+1}-2v_n)\)。代入 Bloch 形式得 \(\begin{pmatrix}2K-M\omega^2&-K(1+e^{-ika})\\-K(1+e^{ika})&2K-m\omega^2\end{pmatrix}(U,V)^T=0\)。非零振型要求行列式为零，矩阵中的相位来自相邻原胞。

(2) 色散行列式化简为 \((2K-M\omega^2)(2K-m\omega^2)-4K^2\cos^2(ka/2)=0\)，亦即 \(Mm\omega^4-2K(M+m)\omega^2+4K^2\sin^2(ka/2)=0\)。解二次方程得 \(\omega_\pm^2=K(1/M+1/m)\pm K\sqrt{(1/M+1/m)^2-4\sin^2(ka/2)/(Mm)}\)。在 \(k=0\)，负支 \(\omega_-(0)=0\)，对应两个原子同相整体平移；正支 \(\omega_+^2(0)=2K(1/M+1/m)\)，对应胞内相对运动，称为光学支。

(3) 在区边界 \(ka=\pi\)，相邻原胞耦合相位使 \(\cos(ka/2)=0\)，动力矩阵的非对角元素在适当相位约定下消失，两个频率恰为 \(\omega_-^2=2K/M\) 与 \(\omega_+^2=2K/m\)。由于 \(M>m\)，声学支最高频与光学支最低频间存在 \((\sqrt{2K/M},\sqrt{2K/m})\) 的无本征模式区间。若 \(M=m\)，两端重合，原胞的人为加倍造成的能带折叠没有真实质量对比造成的频隙。

(4) 在 \(ka\ll1\) 下令 \(S=1/M+1/m\)，将根式展开为 \(S-2\sin^2(ka/2)/(MmS)+\cdots\)，故 \(\omega_-^2\simeq[2K/(M+m)]\sin^2(ka/2)\simeq Ka^2k^2/[2(M+m)]\)。声学长波群速为 \(v_g=d\omega_-/dk\to a\sqrt{K/[2(M+m)]}\)。它随一个原胞的总惯性质量 \(M+m\) 增大而降低，且正比于弹簧刚度平方根；这是离散原子链连续弹性极限的合理量纲与趋势。
